\documentclass[10pt,a4paper]{amsart}
\usepackage[margin=19mm]{geometry}
\usepackage{amsmath,amssymb,mathtools}
\usepackage[scr=rsfs,cal=euler]{mathalfa}
\usepackage{xfrac}
\usepackage[unicode,hidelinks,bookmarks=false]{hyperref}
\newcommand{\ie}{i.e.\ }
\newcommand{\cf}{cf.\ }
\newcommand{\resp}{respectively\ }
\newcommand{\Subsection}{\subsection}
\providecommand{\nobreakdash}{\nobreak\hbox{-}}
\setlength{\emergencystretch}{2em}
\multlinegap0pt
\usepackage{bm}
\usepackage{bbm}
\usepackage{dsfont}
\usepackage{xspace}
\usepackage{cancel}
\usepackage{mathtools}
\usepackage{amsthm}
\makeatletter
\@ifundefined{theorem}{\newtheorem{theorem}{Theorem}}{}
\@ifundefined{lemma}{\newtheorem{lemma}[theorem]{Lemma}}{}
\makeatother
\makeatletter
\renewcommand{\O}{\Omega}
\newcommand{\R}{\mathbb{R}}
\newcommand{\Sw}{\mathcal{S}}
\renewcommand{\d}{\delta}
\newcommand{\g}{\gamma}
\newcommand{\n}{\nabla}
\newcommand{\p}{\partial}
\expandafter\ifx\csname soln\endcsname\relax
\newenvironment{soln}{\begin{proof}[Solution]}{\end{proof}}
\fi
\renewcommand{\t}{\theta}
\newcommand{\vf}{\varphi}
\makeatother
\title[On the Stability of Inverse Conductivity Problem for Polyhed]{On the Stability of Inverse Conductivity Problem for Polyhedral Inclusions under a Single Measurement}
\author{Chun-Hsiang Tsou}
\date{Source: arXiv:2605.17484v2 (CC BY 4.0)}
\begin{document}
\maketitle

\noindent\emph{Source and licence.} On the Stability of Inverse Conductivity Problem for Polyhedral Inclusions under a Single Measurement. Version arXiv:2605.17484v2 (CC BY 4.0), distributed under the Creative Commons Attribution 4.0 International licence, as recorded in the arXiv licence metadata for this exact version and shown on the abstract page https://arxiv.org/abs/2605.17484v2. Full rights notice: licenses/arxiv-2605.17484-CC-BY-4.0.txt.\par\smallskip

\noindent\textit{Editorial scope.} Counted unit: the Lemma whose source label is \texttt{lemma:esti-grad-edge} in the article (source0.tex lines 1294--1300, source label \texttt{lemma:esti-grad-edge}), delivered together with the article's own complete original proof of it (source0.tex lines 1301--1343, one \texttt{proof} environment of 43 lines). No mathematics is written, repaired or re-derived: statement and proof are the article's own text line for line. Required context retained uncounted: eq:edge-singular (lines 1259--1261); eq:embedding-weighted (lines 1281--1283). Every definition, display and label of those lines is retained unchanged. Omitted from this package and not redistributed: 1--1258, 1262--1280, 1284--1293, 1344--1378. Nothing inside a retained excerpt is omitted or altered: every retained body is delivered exactly as the article's lines are written, so no omission receipt of any kind accompanies this package. Citations: the retained excerpts carry no citation command at all, so no bibliography entry of the article is transcribed and no reference is redistributed. Reference adaptation: none was needed and none was made; every reference inside the delivered unit resolves inside it, so no unresolved reference or citation is produced. Printed slips kept literal: no printed word, symbol, hypothesis or number of the counted statement or of its proof was altered. Layout and class: the article's own class is \texttt{amsart} with packages including graphicx, subfigure, amssymb, epstopdf, geometry, bm, amsmath, amsfonts, amsthm, mathrsfs, cite, color, graphics, framed; the wrapper is the lane's pinned \texttt{amsart} editorial wrapper at 10pt on A4, which restates the theorem-like environments the retained text needs and adds the article's own macro definitions that the retained text needs (\texttt{\textbackslash O}, \texttt{\textbackslash R}, \texttt{\textbackslash Sw}, \texttt{\textbackslash d}, \texttt{\textbackslash g}, \texttt{\textbackslash n}, \texttt{\textbackslash p}, \texttt{\textbackslash t}, \texttt{\textbackslash vf}), with extra wrapper packages (bm, bbm, dsfont, xspace, cancel, mathtools). The delivered numbering is the unit's own. Assets: no retained span contains a figure environment, a graphics-inclusion command, a PDF-page inclusion, a table or a TikZ picture; no image file is copied into this package, the package declares no asset dependency and no image file is redistributed with it. Licence: version arXiv:2605.17484v2 of this article is distributed under the Creative Commons Attribution 4.0 International licence, as recorded in the arXiv licence metadata for this exact version and shown on the abstract page https://arxiv.org/abs/2605.17484v2. A full rights notice is delivered at licenses/arxiv-2605.17484-CC-BY-4.0.txt. Characters: every retained excerpt is pure ASCII, so the delivered engine, XeLaTeX with its default Latin Modern text font, reports no missing character.
        \begin{align}\label{eq:edge-singular}
            u_{E,\mu}(r,\t,z)= \mathcal{K}[\g](\widetilde{r},z)\widetilde{r}^\mu\vf(\t).
        \end{align}

        \begin{align}\label{eq:embedding-weighted}
                V_{-s}^{s-\mu}(E)\hookrightarrow V_{-1/2}^{s-\mu}(E) \cong H^{s-\mu}(\R).
        \end{align}

        \begin{lemma}\label{lemma:esti-grad-edge}
            Let $\frac{1}{2}<\mu<s$ and let $u_{E,\mu}$ be given by \eqref{eq:edge-singular} with coefficient $\g\in V_{-s}^{s-\mu}(E)$. Then, in the region $\widetilde{\O}$, there exists a constant $C_{\phi,\mu}>0$ such that
            \begin{align}\label{eq:esti-grad-edge}
                |\n u_{E,\mu}(r,\t,z)| \leq C_{\phi,\mu} \|\g\|_{V_{-s}^{s-\mu}(E)}\frac{r^{\widetilde{\mu}-1}}{\rho^{\widetilde{\mu}}}.
            \end{align}
            Here, $\rho$ denotes the distance from the point $(r,\t,z)$ to its closest endpoint of $E$. Also, $\widetilde{\mu}=\mu$ if $s-\mu\geq 1/2$ and $\widetilde{\mu}=s-1/2$ if $s-\mu<1/2$.
        \end{lemma}

        \begin{proof}
            By a change of variables stretching $E$ into $\R$ and the embedding \eqref{eq:embedding-weighted}, we have $\widetilde{\g}\in H^{s-\mu}(\R)$. We compute $\p_z u_{E,\mu}$ via the chain rule:
            \begin{align*}
                \frac{\p}{\p z}u_{E,\mu}(r,\t,z)&= \left(\frac{\p \widetilde{z}}{\p z}\frac{\p}{\p \widetilde{z}}+\frac{\p \widetilde{r}}{\p z}\frac{\p}{\p \widetilde{r}}\right) u_{E,\mu}(r,\t,z)\\
                &= \frac{1}{\d(z)}\left(\frac{\p}{\p \widetilde{z}}\mathcal{K}[\g]\right)\widetilde{r}^{\mu}\vf(\t)-\frac{\d'(z)}{\d(z)} \left(\widetilde{r}^{\mu+1} \frac{\p}{\p \widetilde{r}}\mathcal{K}[\g]+\mu \widetilde{r}^{\mu}\mathcal{K}[\g]\right) \vf(\t)\\
                \frac{\p}{\p r} u_{E,\mu}(r,\t,z) &= \frac{\p \widetilde{r}}{\p r}\frac{\p}{\p \widetilde{r}} u_{E,\mu}(r,\t,z)=\frac{1}{\d(z)}\left(\frac{\p}{\p \widetilde{r}}\mathcal{K}[\g]\right)\widetilde{r}^{\mu}\vf(\t)+\frac{\mu}{\d(z)} \mathcal{K}[\g]\widetilde{r}^{\mu-1}\vf(\t),\\
                \frac{1}{r} \frac{\p}{\p \t} u_{E,\mu}(r,\t,z) &= \frac{1}{\d(z)}\mathcal{K}[\g]\widetilde{r}^{\mu-1}\vf'(\t).
            \end{align*}
            By the convolution structure of $\mathcal{K}$, we have:
            \begin{align*}
                \frac{\p}{\p \widetilde{z}}\mathcal{K}[\g](\widetilde{r},z)& = \frac{\p}{\p \widetilde{z}}\int_\R \frac{1}{\widetilde{r}}\phi(\frac{\widetilde{z}-t}{\widetilde{r}})\widetilde{\g}(t)dt=\frac{1}{\widetilde{r}}\int_\R \frac{1}{\widetilde{r}}\phi'(\frac{t}{\widetilde{r}})\widetilde{\g}(\widetilde{z}-t) dt,\\
                \frac{\p}{\p \widetilde{r}}\mathcal{K}[\g](\widetilde{r},z)&= -\int_\R \left(\frac{1}{\widetilde{r}^2}\phi(\frac{t}{\widetilde{r}})+\frac{t}{\widetilde{r}^3}\phi'(\frac{t}{\widetilde{r}})\right)\widetilde{\g}(\widetilde{z}-t)dt=-\frac{1}{\widetilde{r}}\int_\R \frac{1}{\widetilde{r}}(t\phi)'(\frac{t}{\widetilde{r}})\g(\widetilde{z}-t)dt.
            \end{align*}
            Note that $\phi',\;(t\phi)'\in \Sw(\R)$. We apply the Sobolev embedding $H^{s-\mu}(\R)\hookrightarrow L^q(\R)$ with $\frac{1}{q}=\frac{1}{2}-(s-\mu)$ if $s-\mu<1/2$ and $q=\infty$ if $s-\mu\geq 1/2$.
            \begin{itemize}
                \item Case $s-\mu <1/2$. Denoting by $q'$ the conjugate exponent of $q$, it follows from H\"older's inequality that
                \begin{align*}
                    \left|\mathcal{K}[\g](\widetilde{r},z)\right| &\leq \widetilde{r}^{-1/q} \|\phi\|_{L^{q'}(\R)} \|\widetilde{\g}\|_{L^q(\R)}=\widetilde{r}^{(s-\mu)-1/2} \|\phi\|_{L^{q'}(\R)} \|\widetilde{\g}\|_{L^q(\R)},\\
                    \left|\frac{\p}{\p \widetilde{z}}\mathcal{K}[\g](\widetilde{r},z)\right| &\leq \widetilde{r}^{(s-\mu)-3/2} \|\phi'\|_{L^{q'}(\R)} \|\widetilde{\g}\|_{L^q(\R)},\\
                    \left|\frac{\p}{\p \widetilde{r}}\mathcal{K}[\g](\widetilde{r},z)\right| &\leq \widetilde{r}^{(s-\mu)-3/2} \|(t\phi)'\|_{L^{q'}(\R)} \|\widetilde{\g}\|_{L^q(\R)}.
                \end{align*}
                \item Case $s-\mu\geq 1/2$. Using Young's inequality, we have
                \begin{align*}
                    |\mathcal{K}[\g](\widetilde{r},z)| &\leq \|\phi\|_{L^1(\R)} \|\widetilde{\g}\|_{L^\infty(\R)},\\
                    \left|\frac{\p}{\p \widetilde{z}}\mathcal{K}[\g](\widetilde{r},z)\right| &\leq \widetilde{r}^{-1} \|\phi'\|_{L^1(\R)} \|\widetilde{\g}\|_{L^\infty(\R)},\\
                    \left|\frac{\p}{\p \widetilde{r}}\mathcal{K}[\g](\widetilde{r},z)\right| &\leq \widetilde{r}^{-1} \|(t\phi)'\|_{L^1(\R)} \|\widetilde{\g}\|_{L^\infty(\R)}.
                \end{align*}
            \end{itemize}
            Define $\widetilde{\mu}:=\mu$ if $s-\mu\geq 1/2$ and $\widetilde{\mu}=s-1/2$ if $s-\mu<1/2$. Substituting these estimates into the expressions for derivatives of $u_{E,\mu}$, we obtain
            \begin{align*}
                \left|\frac{\p}{\p z} u_{E,\mu}(r,\t,z)\right| &\leq \|\widetilde{\g}\|_{L^q(\R)} \left(\frac{\widetilde{r}^{\widetilde{\mu}-1}}{\d(z)} \|\phi'\|_{L^{q'}(\R)} + \frac{|\d'(z)| \widetilde{r}^{\widetilde{\mu}}}{\d(z)}(\|(t\phi)'\|_{L^{q'}(\R)}+\mu \|\phi\|_{L^{q'}(\R)})\right).
            \end{align*}
            The same estimate holds for $\frac{\p}{\p r} u_{E,\mu}$.

            Note that the cut-off function $\eta$ and its support $\ell$ satisfy $\ell \leq 1/2$ and $\eta(\widetilde{r})=0$ if $\widetilde{r}\geq 2\ell$. Hence we may assume $\widetilde{r}\leq 1$. Also, $\d$ is equivalent to the distance to the endpoints of $E$, which implies $\d(z)\geq C_{\d}\rho$ for some constant $C_\d>0$. Then the estimates become
            \begin{align*}
                &\left|\frac{\p}{\p z} u_{E,\mu}(r,\t,z)\right|,\; \left|\frac{\p}{\p r} u_{E,\mu}(r,\t,z)\right|\\
                 &\leq \|\widetilde{\g}\|_{L^q(\R)} C_\d^{-\widetilde{\mu}}\left( \|\phi'\|_{L^{q'}(\R)} + |\d'(z)|(\|(t\phi)'\|_{L^{q'}(\R)}+\mu \|\phi\|_{L^{q'}(\R)})\right)\frac{r^{\widetilde{\mu}-1}}{\rho^{\widetilde{\mu}}},\\
                 & \frac{1}{r} \left|\frac{\p}{\p \t} u_{E,\mu}(r,\t,z)\right| \leq \mu\|\widetilde{\g}\|_{L^q(\R)} C_\d^{-\widetilde{\mu}} \|\phi\|_{L^{q'}(\R)}\frac{r^{\widetilde{\mu}-1}}{\rho^{\widetilde{\mu}}}.
            \end{align*}            
            
            Thus, the claim follows by combining these estimates with \eqref{eq:embedding-weighted}.
        \end{proof}

\end{document}
